\begin{exercisechunk}
\item \textbf{Unobserved responses and probability transfer.}
\label[exercise]{ex:l19-unobserved-responses}
\exerciseprofile{2}{2}{2}{\exproof\quad\excalculation\quad\exinterpretation}{%
Determine which policy probabilities a single comparison leaves unconstrained.}
Let $\cY=\{a,b,c\}$, use a uniform reference and $\beta=1$, and suppose
$\nu(a,b)=1$ and $p(a,b)=3/4$. All policy vectors below list coordinates
in the order $(a,b,c)$.
\begin{enumerate}[(a)]
\item For $t>0$, define $\pi_t=(3,1,t)/(4+t)$. Show that
$q_{\pi_t}(a,b)=3/4$ for every $t$ and that every $\pi_t$ minimizes
the population DPO loss over $\Delta_{\cY}$. Show that
$\pi_t(c)=t/(4+t)$ ranges over $(0,1)$ and explain why the observed
preference probability therefore does not identify the probability of $c$.
\item Let $\pi^{\rm old}=(0.6,0.3,0.1)$ and
$\pi^{\rm new}=(0.3,0.1,0.6)$. Show that
$q_{\pi^{\rm old}}(a,b)=2/3$, $q_{\pi^{\rm new}}(a,b)=3/4$, and
\[
 L_{\rm DPO}(\pi^{\rm old})-L_{\rm DPO}(\pi^{\rm new})
 =\mathrm{kl}(3/4,2/3)>0,
\]
where $\mathrm{kl}$ is binary KL divergence.
Explain how the comparison fit improves while the probabilities of both
$a$ and $b$ decrease, identifying the changed probability ratio and the
response that receives the transferred probability mass.
\end{enumerate}
\begin{hints}
Apply the binary log-loss identity from \cref{prop:population-recovery}
to the single observed pair; no assumption on unobserved preferences is needed.
\end{hints}
\end{exercisechunk}
